\section{Even values, further integral representations, and boundary zeros}
\label{sec:additional}

\subsection{Euler's formula for every positive even integer}
\label{subsec:even-values}
The Bernoulli polynomials of Section~\ref{em:section} and the finite
Fej\'er kernel of Section~\ref{sec:theta-functional} give a short proof
which is independent of the functional equation. Regard $B_{2m}(x)$,
$m\ge1$, as one-periodic. It is continuous across the endpoints and
has circular Lipschitz constant at most $2m C_{2m-1}$ and absolute bound
$C_{2m}$. For every nonzero integer $k$, repeated finite integration by
parts gives
\[
 \int_0^1 B_n(x)e^{-2\pi ikx}\,dx
   =\frac{n}{2\pi ik}\int_0^1B_{n-1}(x)e^{-2\pi ikx}\,dx
       \quad(n\ge2).
\]
The endpoint contribution vanishes because $B_n(1)=B_n(0)$.
At $n=1$ the endpoint difference equals one and the constant derivative
has zero integral against the exponential, so the coefficient is
$-1/(2\pi ik)$. Therefore
\[
 \widehat B_{2m}(k)=-\frac{(2m)!}{(2\pi ik)^{2m}},\qquad
 \widehat B_{2m}(0)=0.
\]
Fej\'er approximation at $x=0$ consequently gives
\[
 B_{2m}=(-1)^{m+1}\frac{2(2m)!}{(2\pi)^{2m}}
       \lim_{N\to\infty}\sum_{k=1}^{N-1}
                  \left(1-\frac{k}{N}\right)k^{-2m}.
\]
This weighted sum has the same named limit as $\zeta(2m)$.
Indeed its difference from that value is at most
\[
 \sum_{k\ge N}k^{-2m}+\frac1N\sum_{k<N}k^{1-2m}
 \le\frac1{N-1}+\frac{1+\log N}{N}\le\frac5{\sqrt N}
       \qquad(N\ge2).
\]
For the last estimate, $\log N\le2\sqrt N$ follows by comparing
$1/x$ with $1/\sqrt x$. A dyadic square $N=2^{2q}$ absorbs this
error with a PR modulus; the Fej\'er error has its own explicit modulus.
We obtain the entire family
\begin{equation}
 \boxed{\displaystyle
 \zeta(2m)=(-1)^{m+1}\frac{B_{2m}(2\pi)^{2m}}{2(2m)!}
                 \qquad(m\ge1).}
 \label{eq:even-values}
\end{equation}
Positivity of the Dirichlet series also proves
$(-1)^{m+1}B_{2m}>0$; no Bernoulli sign theorem was assumed.
In particular
\[
 \zeta(2)=\frac{\pi^2}{6},\quad
 \zeta(4)=\frac{\pi^4}{90},\quad
 \zeta(6)=\frac{\pi^6}{945},\quad
 \zeta(8)=\frac{\pi^8}{9450},\quad
 \zeta(10)=\frac{\pi^{10}}{93555}.
\]
The same Fourier calculation and Fej\'er estimate give the full expansion
\[
 B_{2m}(x)=(-1)^{m+1}\frac{2(2m)!}{(2\pi)^{2m}}
           \sum_{k=1}^{\infty}\frac{\cos(2\pi kx)}{k^{2m}}.
\]
Its tail is at most the prefactor times
$N^{1-2m}/(2m-1)$ after $N$ terms. The case $m=1$, evaluated at any
fixed named $x\in[0,1]$, is a constructive version of the familiar
Fourier-series proof of the Basel evaluation.

The coprime-pair density in \eqref{eq:d-coprime-density} is therefore
$6/\pi^2$. The same number is the squarefree density with an equally
finite proof. Write $Q(N)=\sum_{n\le N}\mu(n)^2$ and
$D=\lfloor\sqrt N\rfloor\ge1$, computed by bounded integer search.
The finite square-divisor identity already proved gives
\[
 Q(N)=\sum_{d\le D}\mu(d)\lfloor N/d^2\rfloor,\qquad
 \left|\frac{Q(N)}N-\frac6{\pi^2}\right|
       \le\frac DN+\frac1D.
\]
The first error counts at most $D$ floor-rounding errors; the second is
$\sum_{d>D}d^{-2}\le1/D$. Choosing $N=2^{2(p+2)}$ gives a PR
modulus smaller than $2^{-p}$.

\subsection{Mellin integral and its endpoint certificates}
\label{subsec:mellin}
For $\sigma=\Re s>1$ with a supplied margin, define
\[
 J(s)=\int_0^\infty\frac{t^{s-1}}{e^t-1}\,dt.
\]
For $0<\epsilon\le1$, the omitted initial interval has absolute integral
at most $\epsilon^{\sigma-1}/(\sigma-1)$, since $e^t-1\ge t$.
For $T\ge1$ and $|s|\le R$, the tail is at most
$2^{R+2}R!e^{-T/2}$, since $(e^t-1)^{-1}\le2e^{-t}$.
On $[\epsilon,T]$, every denominator is explicitly positive and a
rational bound for the derivative follows by the product and quotient
rules. These bounds define $J$ by finite quadrature.

The finite geometric identity gives
$(e^t-1)^{-1}=\sum_{n=1}^Ne^{-nt}+e^{-(N+1)t}/(1-e^{-t})$.
Termwise integration of the finite sum and the substitution $u=nt$
produce $\Gamma(s)\sum_{n\le N}n^{-s}$.
The absolute integral of the omitted positive majorant is at most
\[
 \Gamma(\sigma)\sum_{n>N}n^{-\sigma}
       \le\frac{\Gamma(\sigma)}{\sigma-1}N^{1-\sigma}.
\]
This inequality can itself be obtained first for a finite additional
geometric sum and then from the just-displayed endpoint bounds.
Hence
\begin{equation}
 \boxed{\displaystyle
 \Gamma(s)\zeta(s)=\int_0^\infty\frac{t^{s-1}}{e^t-1}\,dt
                     \quad(\Re s>1).}
 \label{eq:zeta-mellin}
\end{equation}
The gamma lower bound in \eqref{eq:gamma-lower} licenses division when
the conventional quotient version is desired.

There is also a faster-decaying alternating-kernel representation:
\[
 \boxed{\displaystyle
 \Gamma(s)\eta(s)=\int_0^\infty\frac{t^{s-1}}{e^t+1}\,dt
                     \quad(\Re s>0).}
\]
The initial tail is at most $\epsilon^\sigma/\sigma$, and the upper
tail has the same exponential majorant as the gamma integral.
For a finite alternating geometric sum the kernel remainder is
$(-1)^Ne^{-(N+1)t}/(1+e^{-t})$, whose absolute integral is at most
$\Gamma(\sigma)(N+1)^{-\sigma}$. Integrating the finite sum and
using the eta tail therefore proves the formula even when the individual
Dirichlet terms are not absolutely summable. The eta-zeta relation gives
the usual version with factor $(1-2^{1-s})\Gamma(s)$; the cleared
identity is meaningful at its removable points.

One finite integration by parts gives two further formulas:
\begin{align*}
 \Gamma(s+1)\zeta(s)
   &=\int_0^\infty\frac{e^t t^s}{(e^t-1)^2}\,dt
       &&(\Re s>1),\\
 \Gamma(s+1)\eta(s)
   &=\int_0^\infty\frac{e^t t^s}{(e^t+1)^2}\,dt
       &&(\Re s>0).
\end{align*}
At zero the boundary terms are bounded by $\epsilon^{\sigma-1}$
and $\epsilon^\sigma$, respectively; at infinity both have an
exponential bound. These are the exact endpoint conditions needed for
the integrations by parts. The normalizations agree with
\cite[\S25.5]{DLMF25_5}.

\subsection{Fractional-part integral on the right half-plane}
\label{subsec:fractional}
Let $\{x\}=x-\lfloor x\rfloor$, with integration performed separately
on its integer cells. Finite summation by parts, or telescoping the
integrals on those cells, gives for an integer $M\ge1$,
\[
 s\int_1^M\lfloor x\rfloor x^{-s-1}\,dx
       =\sum_{n=1}^M n^{-s}-M^{1-s}.
\]
For $\Re s>1$ this identifies the Dirichlet series with
\begin{equation}
 \zeta(s)=\frac{s}{s-1}
           -s\int_1^\infty\{x\}x^{-s-1}\,dx.
 \label{eq:fractional-zeta}
\end{equation}
The integral itself converges whenever $\sigma>0$, with tail at most
$M^{-\sigma}/\sigma$. The same formula agrees directly with the
Euler--Maclaurin continuation there: write
$\{x\}=\widetilde B_1(x)+1/2$ and integrate once using
$\widetilde B_2'=2\widetilde B_1$. The endpoint contribution is
$-B_2/2$ and the remaining term is
$(s+1)\int_1^\infty\widetilde B_2(x)x^{-s-2}dx/2$.
This is exactly the $m=1$ formula for $H(s)$, including its sign.
Thus \eqref{eq:fractional-zeta} holds for $\sigma>0$, $s\mathrel\#1$,
without another analytic-continuation argument.

If $0<\sigma<1$ with margins at both endpoints, substitute $t=1/x$
and add the explicitly integrable part $t\ge1$, where $\{1/t\}=1/t$.
The rational terms $1/(s-1)$ and $1/(1-s)$ cancel and give
\begin{equation}
 \boxed{\displaystyle
 \int_0^\infty\{1/t\}t^{s-1}\,dt=-\frac{\zeta(s)}s.}
 \label{eq:mathworld-fractional}
\end{equation}
The tail at zero is bounded by $\epsilon^\sigma/\sigma$, and the tail
at infinity by $T^{\sigma-1}/(1-\sigma)$. The finite middle interval
is split at its finitely many reciprocal-integer break points.

\subsection{Unit-cube and logarithmic square integrals}
For every integer $d\ge2$,
\begin{equation}
 \boxed{\displaystyle
 \zeta(d)=\int_{[0,1]^d}
                  \frac{dx_1\cdots dx_d}{1-x_1\cdots x_d}.}
 \label{eq:cube-integral}
\end{equation}
Here is an explicit improper-integral interpretation. Integrate first
on $[0,\rho]^d$ for rational $0<\rho<1$. The geometric series for the
integrand has uniform error at most
$\rho^{d(N+1)}/(1-\rho^d)$ after degree $N$, so finite polynomial
integration gives
\[
 I_\rho=\sum_{j=1}^\infty\frac{\rho^{dj}}{j^d}.
\]
For every $J\ge1$,
\[
 0\le\zeta(d)-I_\rho
 \le dJ(1-\rho)+\frac1{(d-1)J^{d-1}}.
\]
This follows from $1-\rho^{dj}\le dj(1-\rho)$ for $j\le J$, and
from the Dirichlet tail for $j>J$. Taking, for example,
$J=2^{p+4}$ and $1-\rho\le2^{-p-4}/(dJ)$ gives an explicit PR endpoint
modulus, uniformly in $d\ge2$. The finite multiple integrals are just
iterated finite sums, so no measure-theoretic interchange is required.

A further representation appearing on the reference pages is
\begin{equation}
 \boxed{\displaystyle
 \int_0^1\!\int_0^1
       \frac{[-\log(xy)]^s}{1-xy}\,dx\,dy
       =\Gamma(s+2)\zeta(s+2)\quad(\Re s>-1).}
 \label{eq:square-log}
\end{equation}
To specify its endpoints constructively, truncate first to
$\epsilon\le xy\le1-\epsilon$. With $u=xy$, the inner bounds are
$u\le x\le1$ and the Jacobian is $1/x$. Hence its finite double
integral equals
\[
 \int_\epsilon^{1-\epsilon}
            \frac{(-\log u)^{s+1}}{1-u}\,du.
\]
This change of variables takes place on a compact region with explicit
positive denominators and can be checked by two finite sums and the
one-dimensional substitution rule. Now $u=e^{-t}$ changes the integral
to the truncation of
$\int_0^\infty t^{s+1}/(e^t-1)dt$. The endpoint bounds used in
\eqref{eq:zeta-mellin}, with $s$ replaced by $s+2$, are the required
PR modulus. Complex powers of $-\log(xy)>0$ use the positive-real
logarithm convention. For example $s=0$ gives the square case of
\eqref{eq:cube-integral}, and $s=1$ gives $2\zeta(3)$.

\subsection{A quantitative zero-free theorem on $\Re s=1$}
\label{subsec:boundary-zeros}
The Euler logarithm proves the finite-tail form of the classical
$3$--$4$--$1$ inequality:
\[
 1\le\zeta(\sigma)^3|\zeta(\sigma+it)|^4
                       |\zeta(\sigma+2it)|\qquad(\sigma>1).
 \tag{Z1}\label{eq:341-repeat}
\]
For each prime-power angle the relevant numerator is
$3+4\cos u+\cos2u=2(1+\cos u)^2\ge0$. Thus the nonnegativity is
already true for every finite logarithm sum; exponentiation and the
specified tails give (Z1).

Supply integers $A,T\ge1$ with $1/A\le|t|\le T$, and put
\[
 \eta=1/A,\quad R=2T+3,\quad
 B=R(1+\eta^{-1}),\quad M=\eta^{-2}+1+R,
 \quad h=\frac1{256B(M+1)^4}.
 \tag{Z2}\label{eq:boundary-constants}
\]
From \eqref{eq:fractional-zeta}, uniformly for $0\le u\le1$,
\[
 |\zeta(1+u+2it)|\le B,\qquad
 |\zeta'(1+u+it)|\le M.
\]
For the derivative bound differentiate the explicit integral. Its two
integrals have bounds
$\int_1^\infty x^{-\sigma-1}dx=1/\sigma$ and
$\int_1^\infty(\log x)x^{-\sigma-1}dx=1/\sigma^2$.
The derivative of $s/(s-1)$ is $-1/(s-1)^2$. These give exactly $M$;
all differentiated tails have the same explicit logarithmic bounds used
in Section~\ref{em:section}.

Also $\zeta(1+h)\le1+1/h\le2/h$. Equation (Z1) therefore gives
\[
 |\zeta(1+h+it)|^4\ge h^3/(8B)
       =32(M+1)^4h^4,
\]
and hence $|\zeta(1+h+it)|\ge2(M+1)h$. This last inequality follows
by the monotonicity of the fourth power on nonnegative quantities; its
finite rational approximation proof uses a positive rational gap.
Finite integration of the derivative along the horizontal segment gives
$|\zeta(1+h+it)-\zeta(1+it)|\le Mh$. Consequently
\begin{equation}
 \boxed{\displaystyle |\zeta(1+it)|\ge h>0
       \quad\left(1/A\le|t|\le T\right),}
 \label{eq:boundary-apartness}
\end{equation}
with the explicit rational $h$ in (Z2). This is an apartness theorem,
not merely a derivation of $\neg(\zeta(1+it)=0)$.
For an explicit disk, put
\[
 M_* =4\eta^{-2}+2+4R,\qquad
 r=\min\left(\frac12,\frac\eta2,\frac{h}{2M_*}\right).
\]
If $|z-(1+it)|\le r$, then $\Re z\ge1/2$, $|z-1|\ge\eta/2$,
and $|z|\le R$. The differentiated fractional-part formula gives
$|\zeta'(z)|\le M_*$. Segment integration therefore proves
$|\zeta(z)|\ge h/2$ on that entire disk.

The functional equation supplies positive lower bounds on $\Re s=0$
when $|\Im s|$ is bounded above and away from zero. The factors
$\sin(\pi s/2)$, $\Gamma(1-s)$, and $\zeta(1-s)$ each have the
explicit lower bounds already proved. The remaining point is
$\zeta(0)=-1/2$. Together with the reciprocal Dirichlet series and the
left-half-plane sine factor, these give the usual zero-free exterior
statements on their certified charts. The negative even zeros have the
simple derivative formula \eqref{eq:trivial-zero-derivative}.

For a real input $0<x<1$ with rational margins, paired terms show
$\eta(x)\ge1-2^{-x}>0$, while $1-2^{1-x}<0$. Thus $\zeta(x)<0$.
On the negative real axis away from even integers, its sign is the sign
of $\sin(\pi x/2)$ by the functional equation. These are explicit sign
statements, and do not decide equality to a lattice point for an
arbitrary real name.
